% ============================================================
\subsection{Algorithmic Paradigms}
% ============================================================

Each paradigm is identified by a recognition cue in the problem statement.

\begin{itemize}
  \item \textbf{Hashing.} Cue: find a pair, count occurrences, seen before,
        contains a duplicate. Space is traded for $O(1)$ lookup. The canonical
        construction tests whether the complement of the current element
        already appears in a map of previously visited elements.

  \item \textbf{Two pointers.} Cue: a sorted sequence, or a pairing problem.
        Pointers converging from both ends solve ``two elements summing to
        $t$'' in $O(n)$ with constant extra space. Fast and slow pointers
        detect a cycle in a linked list.

  \item \textbf{Sliding window.} Cue: a contiguous subarray or substring
        subject to a constraint. The right edge expands to admit elements and
        the left edge contracts while the constraint is violated. Each element
        enters and leaves once, giving $O(n)$.

  \item \textbf{Prefix sums.} Cue: repeated range-sum queries over a static
        array. Cumulative sums are precomputed once, after which any range
        query is a difference of two entries. This is the array analogue of
        \texttt{SUM() OVER}.

  \item \textbf{Binary search on the answer.} Cue: the smallest $x$ such that
        a condition is feasible. Requires a monotone predicate, meaning that
        feasibility at $x$ implies feasibility at every larger $x$. The search
        ranges over the answer space rather than the input.

  \item \textbf{Sorting as preprocessing.} Cue: intervals, scheduling, or
        grouping equal elements. Sorting converts a global problem into a local
        one in which each element need only be compared against its neighbor.
\end{itemize}

\begin{keybox}[Complexity anchors]
\begin{center}
\begin{tabular}{ll}
\toprule
\textbf{Operation} & \textbf{Cost} \\
\midrule
\texttt{x in set}, \texttt{x in dict}   & $O(1)$ average \\
\texttt{x in list}                      & $O(n)$ \\
\texttt{list.append}                    & $O(1)$ amortized \\
\texttt{list.insert(0,x)}, \texttt{list.pop(0)} & $O(n)$ \\
\texttt{deque.appendleft}, \texttt{deque.popleft} & $O(1)$ \\
\texttt{sorted(xs)}                     & $O(n \log n)$ \\
\texttt{xs[a:b]}                        & $O(b-a)$, copies \\
\bottomrule
\end{tabular}
\end{center}
\end{keybox}
